\documentclass[11pt]{article}
\usepackage[margin=1in]{geometry}
\usepackage[T1]{fontenc}
\usepackage{lmodern,microtype}
\usepackage[colorlinks=true,allcolors=black]{hyperref}
\pagestyle{empty}
\begin{document}
\noindent Mathias Rodr\'iguez Castro\\
Facultad de Ingenier\'ia, Universidad de la Rep\'ublica\\
Montevideo, Uruguay\\
\href{mailto:mathiasr@fing.edu.uy}{mathiasr@fing.edu.uy}

\bigskip
\noindent To the Editors of \emph{Optimization Methods and Software}

\bigskip
\noindent Dear Editors,

\medskip
Please consider ``Pre-Export Row Scaling for Generated Mixed-Integer Programs: Accuracy
Contracts and Software Evaluation'' as a research article in \emph{Optimization Methods
and Software}.

Generated mixed-integer models can have very different numerical representations
without changing their meaning in exact arithmetic. Better matrix conditioning,
however, does not ensure that solver tolerances correspond to the accuracy required
in the original units. The paper addresses this gap through explicit numerical
specifications and independently checked software evaluation.

The work makes two connected contributions. First, matched scaling controls separate
exported conditioning from effects surviving subsequent equilibration and identify
the numerical kernel behind an apparent structural advantage. Second, residual-budget
contracts characterize admissible row factors under original-unit accuracy requirements
and coefficient limits. The construction includes incompatibility, integer-rounding
allowances and exact binary64 export. A reference workflow connects factor selection,
model export, optimization and verification, keeping primal accuracy separate from
the validity of reported mixed-integer bounds.

The evaluation combines hydrothermal dispatch and independently generated
batch-scheduling models, matched scaling controls and enumerated stress tests on
several solvers. A prespecified dispatch study identifies a configuration with lower
work than NumericFocus=3 and no detected contradictions. A paired native-implementation
benchmark reduces complete-process time on a shared host while preserving all 18
paired solver endpoints. Accuracy recovery, adverse outcomes, unresolved bound checks
and parameter sensitivity delimit the findings. Factor admissibility is not a MILP
optimality certificate.

These contributions address the journal's interests in problem scaling, error
sensitivity, validity analysis and optimization software evaluation. Supplementary
Material~1 supplies proofs and protocols; Supplementary Material~2 supplies the
versioned code, data and independent saved-result checks. The public historical
baseline and the current OMS snapshot are identified separately.

I am the sole author and declare no competing interests. No external financial funding
or research grant was received; the historical use of ClusterUY computing resources is
acknowledged in the manuscript.
The manuscript documents AI assistance and my responsibility for its scientific content.
The work has not been published as a journal article and is not under consideration
elsewhere.

\bigskip
\noindent Sincerely,\\[6pt]
Mathias Rodr\'iguez Castro
\end{document}
